%%%%%%%%%%%%%%%%%%%%%%%%%%%%%%%%%%%%%%%%%%%%%%%%%%%%%%%
% % Lines starting with % are comments, which are ignored.
% % This is a handy way of indicating the date and version of
% % your document, to wit:
% %
% % LaTeX sample file
% % Modified March, 2002
% %
%%%%%%%%%%%%%%%%%%%%%%%%%%%%%%%%%%%%%%%%%%%%%%%%%%%%%%%
\DocumentMetadata{
  % PDF standards: ONE key whose value is a comma-separated list --
  % UA-2 (accessibility) plus A-4f (archival; the archival companion of
  % UA-2). Do NOT use separate pdfstandard=... keys: with ordinary
  % key-value parsing only the last one takes effect.
  pdfstandard={UA-2,A-4f},
  lang=en-US,
  tagging=on,
  tagging-setup={
    math/setup=mathml-SE
  }
}

\documentclass{ltx-talk}
% lecture-macros.tex -- shared package setup for lecture-template.tex
% (ltx-talk slides / handouts).
%
% NOTE ON PATHS: every template in this repository is stored one level up
% from where it is compiled (see README.md). All relative \input paths
% below are resolved against the folder the document is compiled in,
% which sits one directory above this file.

% --- Tables, bookmarks, strikethrough ------------------------------------
% booktabs: professional table rules (\toprule, \midrule, \bottomrule)
% bookmark: PDF outline entries (PDF bookmarks)
% ulem[normalem]: \sout strikethrough that survives math mode
\usepackage{booktabs}
\usepackage{bookmark}
\usepackage[normalem]{ulem}
\usepackage{multicol}

% --- Math ------------------------------------------------------------
% amsmath: equation environments and numbering
% mathrsfs: \mathfrak  |  stmaryrd: extra symbols
\usepackage{amsmath}
\usepackage{mathrsfs}
\usepackage{stmaryrd}

% --- Layout -------------------------------------------------------------
% float: [H] float placement
% fontspec: Unicode/OpenType fonts (compile with XeLaTeX or LuaLaTeX)
\usepackage{float}
\usepackage{fontspec}

% --- Colors ---------------------------------------------------------------
% dvipsnames: extended named color palette (the Accessible* colors below
% are built on it);
% table: \rowcolors / \cellcolor for shading tables.
% NOTE: xcolor may only be loaded ONCE per document -- a second
% \usepackage{xcolor} with different options is a fatal "Option clash"
% error, so ALL of its options are listed here.
\usepackage[dvipsnames,table]{xcolor}

% --- Multimedia -----------------------------------------------------------
% animate: embeds audio/video (and legacy Flash) media in a frame.
\usepackage{animate}

% --- Bibliography ----------------------------------------------------------
% biblatex (biber backend): citations for lecture notes.
\usepackage[
  backend=biber,
  style=alphabetic,
  url=true,
  doi=true,
  isbn=true,
  backref=true
]{biblatex}

% --- Colored boxes ----------------------------------------------------------
% talkthemetcolorbox: the colored-box theme for ltx-talk slides
% (shadow/rounded/shaded blocks, showtitle, titlepage styles).
\usepackage[
  shadow,
  rounded,
  shaded,
  blocks,
  showtitle,
  titlepage
]{talkthemetcolorbox}

% --- "theindex" helper environment ------------------------------------------
% Provides a lightweight theindex environment (an itemize) and the
% \@idxitem/\subitem/\subsubitem/\indexspace macros that LaTeX's standard
% theindex machinery expects, so hand-written index pages indent correctly.
\makeatletter
\newenvironment{theindex}
  {\begin{itemize}}
  {\end{itemize}}
\providecommand\@idxitem{\par\hangindent 40\p@}
\providecommand\subitem{\@idxitem \hspace*{20\p@}}
\providecommand\subsubitem{\@idxitem \hspace*{30\p@}}
\providecommand\indexspace{\par \vskip 10\p@ \@plus5\p@ \@minus3\p@\relax}
\makeatother

% --- Index -----------------------------------------------------------------
% imakeidx: builds the printed index.
% The redefinition below appends "|hyperpage" to every index entry so that,
% in the PDF, each index entry is a link to its page.
\usepackage{imakeidx}
\makeatletter
\let\imki@orig@wrindexentrysplit\imki@wrindexentrysplit
\renewcommand\imki@wrindexentrysplit[3]{%
  \imki@orig@wrindexentrysplit{#1}{#2|hyperpage}{#3}%
}

\makeatletter
% --- Section titles: capture + one bookmark per logical section -------------
% \section is redefined so that:
%   1. the section title is captured into the plain-TeX register
%      \mytalksectitletoks (reprinted verbatim on the auto divider frame,
%      see the cmd/section/after hook below -- plain-TeX storage is
%      catcode-safe, unlike expanding the title in a hook), and
%   2. exactly ONE depth-1 PDF bookmark is created per \section call
%      (hyperref's own bookmark is suppressed via \Hy@writebookmark).
\DeclareCommandCopy{\OriginalSectionForBookmarks}{\section}
\newtoks\mytalksectitletoks

\RenewDocumentCommand{\section}{s o m}{%
  \begingroup
    \def\Hy@writebookmark##1##2##3##4##5{}%
    \global\mytalksectitletoks={#3}%
    \IfBooleanTF{#1}
      {\phantomsection\IfNoValueTF{#2}{\OriginalSectionForBookmarks*{#3}}{\OriginalSectionForBookmarks*[#2]{#3}}}%
      {\IfNoValueTF{#2}{\OriginalSectionForBookmarks{#3}}{\OriginalSectionForBookmarks[#2]{#3}}}%
  \endgroup
  \IfBooleanTF{#1}
    {\IfNoValueTF{#2}{\pdfbookmark[1]{#3}{section-star.\thepage.\number\inputlineno}}{\pdfbookmark[1]{#2}{section-star.\thepage.\number\inputlineno}}}%
    {\IfNoValueTF{#2}{\pdfbookmark[1]{Section \thesection: #3}{section.\thesection}}{\pdfbookmark[1]{Section \thesection: #2}{section.\thesection}}}%
}
\makeatother

% --- One PDF bookmark per logical frame (NOT per overlay) ---------------------
% ltx-talk writes one PDF page per overlay, so a naive frametitle hook would
% register a new bookmark for every overlay of the same frame. The
% \g_slideutils_frame_bookmark_int register remembers the last frame that
% got a bookmark, so subsequent overlays of the same frame are skipped.
% This is also what keeps BeamerPresenter's cumulative annotations working:
% one logical frame must stay one logical unit in the PDF outline.
\ExplSyntaxOn
\int_new:N \g_slideutils_frame_bookmark_int

\AddToHook{cmd/frametitle/after}{%
  \tl_if_empty:NF \g__talk_frame_title_tl
    {
      \ifnum\value{frame}=\g_slideutils_frame_bookmark_int\relax
      \else
        \int_gset:Nn \g_slideutils_frame_bookmark_int { \value{frame} }%
        \pdfbookmark[2]{\tl_use:N \g__talk_frame_title_tl}
          {frame.\int_use:N \g_slideutils_frame_bookmark_int}
      \fi
    }
}
\ExplSyntaxOff
\makeatother

% --- Printed index + theme for the colored boxes -------------------------------
% \makeindex[...]: two-column index titled "Index", listed in the TOC.
% The two EditInstance calls style the standard header/footer instances
% of talkthemetcolorbox (cyan background; footer order title/author/
% date/framenumber).
\makeindex[columns=2, title=Index, intoc]\EditInstance{header}{std}{
  background-color = cyan,
  color            = black,
  height           = 0.75cm
}
\EditInstance{footer}{std}{
  background-color = cyan,
  element-order = {title, author, date, framenumber},
}

% --- Auto section divider frame ---------------------------------------------------
% After each \section, insert a centered, large, structure-colored frame
% that prints the section title (from \mytalksectitletoks).
\AddToHook{cmd/section/after}{%
  \begin{frame}
    \centering
    \color{structure}
    \LARGE
    \the\mytalksectitletoks
  \end{frame}%
}

% --- Tagging sockets (tagged PDF structure) ---------------------------------------
% ltx-talk exposes tagging "sockets" where the PDF structure can attach
% section titles; register one and leave it unassigned (none) so the
% class's own tagging keeps control.
\ExplSyntaxOn
\AtBeginDocument{
  \NewTaggingSocketPlug { talk / sec / title } { none } { }
  \AssignTaggingSocketPlug { talk / sec / title } { none }
}
\ExplSyntaxOff

% --- Accessibility roles ----------------------------------------------------------------
% Map \frametitle to the H1 heading role in the tagged PDF, so screen
% readers treat slide titles as headings.
\tagpdfsetup{role / new-tag = frametitle / H1}

\makeatletter
% --- Page label = frame number -------------------------------------------------------
% Redefine what feeds the PDF page label to track the frame counter,
% not the physical page counter. Replace \theframenumber below with
% whatever internal macro/property ltx-talk actually uses to produce
% the "framenumber" footer element — confirm via ltx-talk-frame.dtx.
\renewcommand\thepage{\arabic{frame}}
\makeatother

\addbibresource{refs.bib}

% --- "Problem Time!" slides ------------------------------------------------------------------
% \ProblemSlide{problem}{solution} emits two frames: the problem, then
% its solution after a \pause, each numbered by the problem counter.
\newcounter{problemcounter}
\setcounter{problemcounter}{0}
\newcommand{\ProblemSlide}[2]{
  \stepcounter{problemcounter}
  \begin{frame}
    \frametitle{Problem Time!}
    Problem \theproblemcounter: #1
  \end{frame}

  \begin{frame}
    \frametitle{Problem Time!}
    Solution to Problem \theproblemcounter: \pause

    #2
  \end{frame}
}
% \highlight[<colour>]{<stuff>}
% from https://tex.stackexchange.com/questions/33401/a-version-of-colorbox-that-works-inside-math-environments
\newcommand{\highlight}[2][yellow]{\mathchoice%
  {\colorbox{#1}{$\displaystyle#2$}}%
  {\colorbox{#1}{$\textstyle#2$}}%
  {\colorbox{#1}{$\scriptstyle#2$}}%
  {\colorbox{#1}{$\scriptscriptstyle#2$}}}%

% --- Fonts ------------------------------------------------------------------------------
% Lato (text) + LeteSansMath (math): open, accessible typefaces.
\setmainfont{Lato}[
  Ligatures=TeX,
  Numbers={Lining,Monospaced}
]
\setsansfont{Lato}[
  Ligatures=TeX,
  Numbers={Lining,Monospaced}
]
\setmathfont{LeteSansMath.otf}\usepackage{emoji}

% --- Accessible colors ---------------------------------------------------------------------
% WCAG-contrast-safe variants of red/green for text on white:
% red!80!black and green!40!black both clear the 4.5:1 requirement.
% AltYellow: soft background for highlight boxes.
\colorlet{AccessibleRed}{red!80!black}
\colorlet{AccessibleGreen}{green!40!black}
\definecolor{AltYellow}{RGB}{255,245,170}

% --- Links ---------------------------------------------------------------------------------
% hyperref + link colors: every link uses AccessibleGreen (7.2:1 on
% white), so internal, citation, and URL links all pass WCAG 1.4.3.
\usepackage{hyperref}
\hypersetup{
  colorlinks=true,
  linkcolor=AccessibleGreen,
  urlcolor=AccessibleGreen,
  citecolor=AccessibleGreen,
}

% --- Shared course macros ---------------------------------------------------------------------
% Root macros.tex: math shorthands used across all course documents
% (\RR, \NN, \ZZ, \QQ, \grad, \curl, \tr, \der, \pder, ...).
\newcommand{\RR}{\mathbb{R}}
\newcommand{\NN}{\mathbb{N}}
\newcommand{\ZZ}{\mathbb{Z}}
\newcommand{\QQ}{\mathbb{Q}}
\newcommand{\vv}{\vec{v}}
\newcommand{\ww}{\vec{w}}
\newcommand{\xx}{\vec{x}}
\newcommand{\yy}{\vec{y}}
\newcommand{\zz}{\vec{z}}
\newcommand{\uu}{\vec{u}}
\newcommand{\basis}[1]{\vec{e_{#1}}}
\newcommand{\id}{\mathrm{id}}
\newcommand{\CC}{\mathbb{C}}
\renewcommand{\div}{\operatorname{div}}
\newcommand{\grad}{\operatorname{grad}}
\newcommand{\curl}{\operatorname{curl}}
\newcommand{\tr}{\operatorname{tr}}
\newcommand{\der}[2]{\dfrac{d #1}{d #2}}
\newcommand{\pder}[3][]{\dfrac{\partial^{#1} #2}{\partial #3}}


% % Title and author(s)
%%%%%%%%%%%%%%%%%%%%%%%%%%%%%%%%%%%%%%%%%%%%%%%%%%%%%%%
\title{Lecture 1: Class Overview and Background Review}
\author{Garrett Credi}
\date{09/08/2026}
%%%%%%%%%%%%%%%%%%%%%%%%%%%%%%%%%%%%%%%%%%%%%%%%%%%%%%%
% %
% % The next command allows your in import encapsulated
% % postscript files, .epsf or .eps files, which
% % contain vector graphic image data.
% %
%%%%%%%%%%%%%%%%%%%%%%%%%%%%%%%%%%%%%%%%%%%%%%%%%%%%%%%
% %   The first thanks indicates your affiliation
% %
% %  Just the name here.
% %
% % Your mailing address goes at the end.
% %
% % \thanks is also how you indicate grant support
% %
%%%%%%%%%%%%%%%%%%%%%%%%%%%%%%%%%%%%%%%%%%%%%%%%%%%%%%%

%%%%%%%%%%%%%%%%%%%%%%%%%%%%%%%%%%%%%%%%%%%%%%%%%%%%%%%
% $ The next command ensures the margins are 1 inch on
% % all sides.
% %
%%%%%%%%%%%%%%%%%%%%%%%%%%%%%%%%%%%%%%%%%%%%%%%%%%%%%%%
\begin{document}
\begin{frame}
  \maketitle
\end{frame}

\begin{frame}
  \frametitle{Welcome to MATH 1271!}

  \begin{figure}[htpb]
    \centering
    \includegraphics[height=0.7\textheight]{1271Banner.jpg}
    \caption{Our Mascot, Maxwell, Deep in Thought}
  \end{figure}
\end{frame}

\begin{frame}
  \frametitle{Today's Goals}
  \begin{enumerate}
    \item Overview of the Course \pause
    \item Lecture on Algebra/Precalc Background \pause
    \item Worksheet Time!
  \end{enumerate}
\end{frame}

\section{Course Overview}

\begin{frame}
  \frametitle{Course Policies I}

  I've structured this course in the following manner: \pause
  \begin{enumerate}
    \item We will start class with a \textit{5 minute} quiz on the topic(s) covered in the previous lecture \pause
    \item We will then begin an ~75 minute lecture, with a ~5-ish minute break about halfway through \pause
    \item After another short break I'll hand out worksheets where we'll get a chance to apply the concepts we learned during lecture
   \end{enumerate}
\end{frame}

\begin{frame}
  \frametitle{Course Policies II}

  After each lecture you will have some \textit{homework} due at 11:59 PM the following day. \pause
  It will:
  \begin{enumerate}
    \item Be graded on an \textit{honest completion} basis \pause
    \item Be split up into certain ``levels" \pause
    \item And will conclude with one ``challenge problem" designed to help you all develop more confidence for quiz/exam environments
  \end{enumerate}
\end{frame}

\begin{frame}
  \frametitle{Course Policies III}

  We will also have \textit{weekly engagement portfolios} \pause
  which are short ~100 word pieces of writing to get you to engage in some metacognition about
  how you are learning in this course. \pause 

  These will be consistently due at 11:59 PM on Mondays.
\end{frame}
\begin{frame}
  \frametitle{Course Policies IV}

  As mentioned before, we will have quizzes at the start of class \pause
  relating to the content included in the previously-due homework, graded
  for accuracy.

  We will also have \textit{three} in class exams: \pause
  \begin{enumerate}
    \item Exam 1: 10/6 \pause
    \item Exam 2: 11/3 \pause
    \item Exam 3: 12/8 \pause
  \end{enumerate}
  Which will all take place during our regular class meeting here (Tate Hall 120). \pause

  Finally we will have a \textit{cumulative} final exam on 12/17 from 12PM-3PM Location TBD.
\end{frame}

\begin{frame}
  \frametitle{Course Policies V}

  Some things that might make you happy! \pause
  \includegraphics[height=0.3\textheight]{yippee.jpg} \pause

  \begin{itemize}
    \item (At this point) your two lowest quizzes and two lowest engagement portfolios will be dropped. \pause
    \item If you do poorly on an exam, you can have your performance on the final exam replace the grade if it is better. \pause
    \item All this information is easily accessible and searchable on our syllabus on Canvas! \pause
  \end{itemize}
\end{frame}

\begin{frame}
  \frametitle{Office Hours?}

  I'm planning on holding two hours of office hours a week: \pause
  in order that I hold them at the most opportune times please indicate 
  your desires/availability on the \href{https://www.when2meet.com/?38405370-tAcFp}{When2Meet here}.

  There are also Office Hours that are held by some other TAs that you can find
  \href{https://cse.umn.edu/math/tutoring-and-math-lab-office-hours}{here.}
\end{frame}

\begin{frame}
  \frametitle{Who Am I?}

  Hi! I'm Garrett! I am: \pause
  \begin{enumerate}
    \item A 3rd Year PhD student in the math department \pause
      \begin{enumerate}
        \item My research is in a field of math called \textit{algebraic topology} \pause
      \end{enumerate}
    \item Someone who has TA'd for math courses three times (1x1271, 2x3283W) and taught once (1x2263) \pause
    \item The proud cat dad of the lovely, gracious, cute, and venerable Maxwell (who is our course banner on canvas!) \pause
    \item Super excited to teach you all some calculus this semester!!! \pause
  \end{enumerate}

  You can contact me at \href{mailto:credi008@umn.edu?subject=[MATH 1271]}{\nolinkurl{credi008@umn.edu}}
\end{frame}

\begin{frame}
  \frametitle{Who are you?}

  What is your:
  \begin{enumerate}
    \item Name \pause
    \item Major and standing (freshman, sophomore, etc)
    \item Favorite \textit{thing}?
  \end{enumerate}
\end{frame}

\section{Algebra and Precalculus Review}

\begin{frame}
  \frametitle{Let's Do Some Math!}

  I want to spend today making sure we're all clear
  on the mathematical background
  necessary for success in this class: \pause
  \begin{enumerate}
    \item Functions, their equations, graphs, and properties (e.g. domain and range)\pause
    \item Techniques with fractions \pause
    \item Trigonometric functions \pause
    \item Solving Equations
  \end{enumerate}
\end{frame}

\begin{frame}
  \frametitle{Functions I}

  What is a function? \pause

  \begin{block}{Function}
    A \index{function}function is an association from ``things" of type $A$ to
    other ``things" of type $B$, typically denoted as $f : A \rightarrow B$. \pause

    For the purposes of our class, we will study \textbf{real-valued}\index{function!real-valued} functions
    of a \textbf{real variable} \pause e.g. functions
    that take in a real number $x \in \RR$ (read as $x$ is in the real numbers)
    and output a real number $y \in \RR$. \pause
  \end{block}

  \begin{block}{Equation of a Function}
    Real valued functions of a single real variable (we'll just call these functions from now on)
    can be \textbf{defined} by an equation \index{function!equation} \pause
    which is just an algebraic specification of what the function does to an input number. \pause
  \end{block}

  E.g. we can define the \textit{squaring} function as $f(x) = x^2$. \pause
  This means the value of the function $f$ at some real number $x$ is the square of $x$,
  so, for instance, $f(2) = 2^2 = 4, f(3) = 3^2 = 9,$ etc.
\end{frame}

\begin{frame}
  \frametitle{Functions II}

  \begin{block}{Graph of a function}
    The \textbf{graph}\index{function!graph} of a function $f : \RR \rightarrow \RR$ \pause
    is a graphical/geometric way of interpreting what a function does. \pause

    To draw the graph, we first draw a \textit{2D-Coordinate Grid}, \pause
    and then draw all the points $(x, f(x))$ connected via a smooth line.
  \end{block}

  E.g. \vspace{\stretch{1}}
\end{frame}

\begin{frame}
  \frametitle{Functions III}

  Sometimes functions behave differently depending on what kind of real number we input into them \pause
  for instance the \textit{absolute value function}\index{absolute value} \pause
  leaves positive real numbers unchanged ($|2| = 2$)
  but turns negative numbers into their opposites ($|-2| = 2$). \pause
  How can we define these functions?

  \begin{block}{Piecewise Function}
    A function is \textbf{defined piecewise}\index{function!piecewise} by
    the following method: \pause
    \begin{equation}
      f(x) = \begin{cases} f_1(x), & \text{some condition on x} \\
        f_2(x), & \text{some other condition on x} \\
        f_3(x), & \text{maybe some other condition on x} \\
        \vdots \\
        g(x), & \text{otherwise if none of the above conditions are met}
      \end{cases}
    \end{equation} \pause
  \end{block}
\end{frame}
\begin{frame}
  \frametitle{Functions IV}
  For example, we define the absolute value function as: \pause
  \begin{equation}
    |x| = \begin{cases} x, & x \geq 0 \\ -x, & x \leq 0 \end{cases}
  \end{equation} \pause

  and we can define some wacky function as: \pause
  \begin{equation}
    wacky(x) = \begin{cases} x^2, & x > -5 \\ -x^3 + 2x + 1, & -5 \leq x < 2 \\ 1, & 2 \leq x < 3 \\ \dfrac{x}{2} + 3, & \text{otherwise} \end{cases} 
  \end{equation}
\end{frame}

\begin{frame}
  \frametitle{Functions V}
  There are two important properties of a function that we need to talk about: domain and range. \pause

  \begin{block}{Domain and Range of a Function}
    The domain of a function \index{function!domain} is the set of all values
    the function accepts as input. \pause
    It is kind of arbitrary! \pause
    However, we'll follow the convention that the 
    domain should be the \textbf{largest} set of values
    it could accept as input. \pause

    Similarly, the range of a function \index{function!range}
    is the set of all values the function could possibly output
    over all possible inputs. Given a specified domain, 
    there is \textbf{exactly one} set of ranges \pause
    so we don't have to deal with any arbitrariness.
  \end{block}

  For example, the \textit{square-root} function $f(x) = \sqrt{x}$
  would have:
  \begin{itemize}
    \item Domain: \pause the \textit{non-negative real numbers} $[0, \infty)$ \pause
    \item Range: \pause again the \textit{non-negative real numbers} $[0, \infty)$ \pause
  \end{itemize}

  What about a function like $f(x) = x^2 + 1$?
\end{frame}

\begin{frame}
  \frametitle{Fractions I}

  It's also worth going over some fraction foundations which will be
  very useful very soon in the course:

  \begin{block}{Addition of Fractions}
    \vspace{0.5in}
  \end{block} \pause

  \begin{block}{Multiplication of Fractions}
    \vspace{0.5in}
  \end{block}

\end{frame}

\begin{frame}
  \frametitle{Fractions II}
  I also want to make clear what people mean by: \pause
  \begin{itemize}
    \item Clearing out common factors/
    \item Factoring [blah] out/
    \item Multiplying by $1$
  \end{itemize}

  \begin{alertblock}{The Multiplying by $1$ Trick}
    \vspace{1in}
  \end{alertblock}
\end{frame}

\begin{frame}
  \frametitle{Special Functions}

  We have a class of ``special" functions: \pause
  \begin{itemize}
    \item The exponential function $\exp(x) = e^x$ \pause
    \item It's inverse, the natural logarithm $\ln(x)$ \pause
    \item The power functions $x^n$ \pause
    \item Their inverses $\sqrt[n]{x}$ \pause
    \item The trigonometric functions 
  \end{itemize}
\end{frame}

\begin{frame}
  \frametitle{Trigonometric Functions}
\end{frame}

\begin{frame}
  \frametitle{Solving Equations I}
  Oftentimes we're interested in finding points where
  functions attain certain values or exhibit important behavior, \pause
  how do we find them? \pause

  For instance, let's say we want to buy some ramen
  and the cost of ramen in dollars varies by weight in pounds according to the
  following function: \pause
  $$P(w) = x^2 + 7x.$$ \pause

  If we only have 30 dollars, how much ramen can we buy?
\end{frame}

\begin{frame}
  \frametitle{Solving Equations II}
  How do we set that up mathematically? \pause

  We want to find the weight $w$ so that:
  \begin{equation}
    P(w) = \$30  
  \end{equation} \pause

  So we can use our definition of $P$ to get: \pause

  \begin{gather*}
    w^2 + 7w = 30 \pause \\
    w^2 + 7w - 30 = 0 \pause \\
    (w + 10)(w - 3) = 0 \pause \\
    w = -10 \text{ or } w = 3
  \end{gather*}

  And since $-10$ lbs of ramen doesn't make physical sense,
  we get that we can by $3$ lbs of ramen with that much money.
\end{frame}

\begin{frame}
  \frametitle{Solving Equations III}
  Why did we choose to solve
  $$(w+10)(w-3) = 0$$
  for $w = -10, w = 3$

  instead of solving
  $$w(w+7)=30$$
  for $w = 30$ or $w + 7 = 30$? \pause

  \begin{alertblock}{No Zero Divisors}
    If there are two real numbers $a,b$ so that $ab = 0$, \pause
    it \textbf{MUST} be the case that $a = 0$ or $b = 0$. \pause
  \end{alertblock}

  Question: If there are two real numbers $a,b$ so that $ab = 30$,
  must it be the case that $a = 30$ or $b = 30$?
\end{frame}

\begin{frame}
  \frametitle{Solving Equations IV}

  How might we approach solving for $x$ in an equation like: \pause

  \begin{equation}
    (e^x)^2 - 5e^x + 6 = 0?
  \end{equation} \pause

  Or maybe like: \pause

  \begin{equation}
    2\sin(x)^2 + \cos(x)^2 = 2?
  \end{equation}
\end{frame}

\begin{frame}
  \frametitle{Worksheet Time}

  Let's get into groups and work on the worksheet!
\end{frame}

\begin{frame}
  \frametitle{End Of Class}

  That's all folks!
\end{frame}
\end{document}
